\section{Motivasyon}
\label{IkinciBolum}

Türk İşaret Dili, işitme engelli bireylerin hem kendi aralarında hem de
çevreleriyle iletişim kurarken kullandığı doğal dildir~\cite{TIF2023}. Ancak
işaret dilini bilmeyen kişilerle kurulan iletişim çoğu zaman yazı yazmaya ya da
bir yakının tercümanlığına bağlı kalmaktadır. Hastane, banka, okul ve resmî
daire gibi ortamlarda bu durum hem zaman kaybına hem de ciddi bir mahremiyet
sorununa yol açmaktadır.

Projenin çıkış noktası, bu iletişim barındaki en yaygın cihazı kullanmaktır.
Eldiven, derinlik kamerası gibi ek donanım gerektirmeyen bir çözüm, kullanıcı
açısından kabul edilebilirliği artırır ve maliyeti düşürür. Bu nedenle sistem
yalnızca telefon kamerası ve cihaz üzerinde çalışan bir model ile kurgulanmıştır.
İnternet bağlantısına ihtiyaç duymaması, uygulamanın sınıfta, hastanede veya
şebekenin zayıf olduğu yerlerde de kullanılabilmesini sağlar.

Beklenen katkı, kelime düzeyinde çalışan bir tanıma modülünün mobil cihazda
gerçek zamanlıya yakın hızda çalıştığını göstermektir. Cümle çevirisi projenin
kapsamı dışında bırakılmış ve gelecek çalışma olarak not edilmiştir.

Çalışma üç sorumluluk alanına ayrılmıştır: mobil uygulama, kullanıcı arayüzü ve
donanım entegrasyonu (Ahmet Çelik); yapay zekâ modeli, eğitim ve optimizasyon
(Mustafa Göktepe); veri hazırlığı, veri hattı ve normalizasyon (Mustafa Küpeli).
Bu raporun yöntem bölümü ağırlıklı olarak model tarafını anlatmaktadır.
